% !TEX program = lualatex
% Transparencias de Fundamentos de las Matemáticas II
% Clase 03: Divisibilidad I

\documentclass[aspectratio=169,11pt]{beamer}

\usepackage{fontspec}
\usepackage[spanish,es-nodecimaldot]{babel}
\usepackage{amsmath,amssymb,mathtools}
\usepackage{booktabs}
\usepackage{csquotes}
\usepackage{xcolor}

\usetheme{Madrid}
\usecolortheme{default}
\setbeamertemplate{navigation symbols}{}
\setbeamertemplate{footline}[frame number]

\definecolor{FdMBlue}{RGB}{20,55,100}
\definecolor{FdMRed}{RGB}{130,30,30}
\definecolor{FdMGreen}{RGB}{25,100,60}

\setbeamercolor{structure}{fg=FdMBlue}
\setbeamercolor{title}{fg=white,bg=FdMBlue}
\setbeamercolor{frametitle}{fg=white,bg=FdMBlue}
\setbeamercolor{block title}{fg=white,bg=FdMBlue}
\setbeamercolor{block body}{bg=blue!3}
\setbeamercolor{alerted text}{fg=FdMRed}

\setmainfont{Latin Modern Roman}
\setsansfont{Latin Modern Sans}
\setmonofont{Latin Modern Mono}

\newcommand{\N}{\mathbb{N}}
\newcommand{\Z}{\mathbb{Z}}
\newcommand{\Q}{\mathbb{Q}}
\newcommand{\R}{\mathbb{R}}
\newcommand{\C}{\mathbb{C}}
\newcommand{\Pow}{\mathcal{P}}
\newcommand{\id}{\operatorname{id}}
\newcommand{\mcd}{\operatorname{mcd}}
\newcommand{\mcm}{\operatorname{mcm}}
\newcommand{\im}{\operatorname{Im}}

\title[FdM II -- Clase 03]{Divisibilidad I}
\subtitle{Divisores, múltiplos y números primos}
\author{Antonio Falcó Montesinos}
\institute{Universidad CEU Cardenal Herrera}
\date{Fundamentos de las Matemáticas II}

\begin{document}

\begin{frame}
  \titlepage
\end{frame}

\begin{frame}{Objetivos de la clase}
\begin{itemize}
  \item Definir la relación \(a\mid b\).
  \item Trabajar con divisores y múltiplos.
  \item Introducir números primos y compuestos.
\end{itemize}
\end{frame}

\begin{frame}{Mapa de la clase}
\tableofcontents
\end{frame}

\section{Divisibilidad}

\begin{frame}{Divisibilidad}
\begin{block}{Definición}
\begin{itemize}
  \item Sean \(a,b\in\mathbb{Z}\).
  \item Diremos que \(a\mid b\) si existe \(q\in\mathbb{Z}\) tal que \(b=aq\).
  \item La definición debe expresarse mediante una existencia.
\end{itemize}
\end{block}
\end{frame}

\begin{frame}{Divisibilidad}
\begin{block}{Ejemplos}
\begin{itemize}
  \item \(3\mid 12\), porque \(12=3\cdot 4\).
  \item \(5\nmid 12\), porque no existe entero \(q\) con \(12=5q\).
  \item Todo entero no nulo divide a \(0\).
\end{itemize}
\end{block}
\end{frame}

\begin{frame}{Divisibilidad}
\begin{block}{Propiedades}
\begin{itemize}
  \item Si \(a\mid b\) y \(a\mid c\), entonces \(a\mid (mb+nc)\) para \(m,n\in\mathbb{Z}\).
  \item Si \(a\mid b\) y \(b\mid c\), entonces \(a\mid c\).
  \item La divisibilidad traduce multiplicatividad entera.
\end{itemize}
\end{block}
\end{frame}

\section{Primos}

\begin{frame}{Primos}
\begin{block}{Número primo}
\begin{itemize}
  \item Un primo es un entero positivo \(p\geq 2\) cuyos únicos divisores positivos son \(1\) y \(p\).
  \item El número \(1\) no se considera primo.
  \item Esta convención preserva la unicidad de la factorización.
\end{itemize}
\end{block}
\end{frame}

\begin{frame}{Primos}
\begin{block}{Compuesto}
\begin{itemize}
  \item Un entero \(n\geq 2\) es compuesto si no es primo.
  \item Equivalentemente, \(n=ab\) con \(1<a<n\) y \(1<b<n\).
  \item Ejemplo: \(12=3\cdot 4\).
\end{itemize}
\end{block}
\end{frame}

\section{Profundización}

\begin{frame}{Profundización}
\begin{block}{Error habitual}
\begin{itemize}
  \item \(a\mid b\) no significa que \(a/b\) sea una operación que se esté realizando.
  \item Significa que \(b\) es múltiplo entero de \(a\).
  \item La variable auxiliar \(q\) debe pertenecer a \(\mathbb{Z}\).
\end{itemize}
\end{block}
\end{frame}

\begin{frame}{Profundización}
\begin{block}{Cero}
\begin{itemize}
  \item Si \(a\neq 0\), entonces \(a\mid 0\).
  \item La afirmación \(0\mid b\) solo es cierta cuando \(b=0\).
  \item Por eso muchos enunciados excluyen divisores nulos.
\end{itemize}
\end{block}
\end{frame}

\begin{frame}{Profundización}
\begin{block}{Práctica}
\begin{itemize}
  \item Listar los divisores positivos de \(18\).
  \item Decidir si \(7\mid 91\).
  \item Justificar por qué \(1\) no debe ser primo.
\end{itemize}
\end{block}
\end{frame}


\section{Detalle y práctica}

\begin{frame}{Detalle y práctica}
\begin{block}{Divisibilidad como relación}
\begin{itemize}
  \item Escribimos \(a\mid b\) cuando existe \(k\in\mathbb{Z}\) tal que \(b=ak\).
  \item La relación es reflexiva en \(\mathbb{Z}\setminus\{0\}\): \(a\mid a\).
  \item Es transitiva: si \(a\mid b\) y \(b\mid c\), entonces \(a\mid c\).
\end{itemize}
\end{block}
\end{frame}

\begin{frame}{Detalle y práctica}
\begin{block}{Actividad breve}
\begin{itemize}
  \item Decidir si \(6\mid 42\), \(8\mid 42\) y \(-5\mid 35\).
  \item Probar que si \(a\mid b\) y \(a\mid c\), entonces \(a\mid (b+c)\).
  \item Indicar por qué \(0\mid b\) solo puede ocurrir si \(b=0\).
\end{itemize}
\end{block}
\end{frame}

\begin{frame}{Cierre}
\begin{itemize}
  \item La divisibilidad se define con una existencia entera.
  \item Los primos son los bloques básicos de la multiplicación.
  \item El caso del cero debe tratarse con cuidado.
\end{itemize}
\end{frame}

\begin{frame}{Trabajo recomendado}
\begin{itemize}
  \item Releer las secciones correspondientes del manual.
  \item Identificar definiciones, símbolos nuevos y enunciados principales.
  \item Preparar dudas y ejemplos para el seminario asociado.
\end{itemize}
\end{frame}

\end{document}
